% Compile with LuaLaTeX or XeLaTeX. The working source is Ukrainian; the same
% structure can be retained when the article is translated into English.
\usepackage{iftex}
\ifPDFTeX
  \PackageError{dudyk-manuscript}{Compile this manuscript with LuaLaTeX or XeLaTeX}{}
\fi

\usepackage{fontspec}
\usepackage{polyglossia}
\setdefaultlanguage{ukrainian}
\setotherlanguage{english}
% TeX Gyre Termes can be present without Cyrillic glyphs. Prefer the
% Windows fonts, which cover Ukrainian, and use DejaVu elsewhere.
\IfFontExistsTF{Times New Roman}{
  \setmainfont{Times New Roman}
  \setsansfont{Arial}
  \setmonofont{Consolas}
  \newfontfamily\cyrillicfont{Times New Roman}[Script=Cyrillic]
  \newfontfamily\cyrillicfontsf{Arial}[Script=Cyrillic]
  \newfontfamily\cyrillicfonttt{Consolas}[Script=Cyrillic]
}{
  \setmainfont{DejaVu Serif}
  \setsansfont{DejaVu Sans}
  \setmonofont{DejaVu Sans Mono}
  \newfontfamily\cyrillicfont{DejaVu Serif}[Script=Cyrillic]
  \newfontfamily\cyrillicfontsf{DejaVu Sans}[Script=Cyrillic]
  \newfontfamily\cyrillicfonttt{DejaVu Sans Mono}[Script=Cyrillic]
}

\usepackage[a4paper,margin=25mm]{geometry}
\usepackage{microtype}
\usepackage{setspace}
\setstretch{1.08}

\usepackage{amsmath,amssymb,mathtools,bm}

\usepackage{graphicx}
\graphicspath{{../../figures/}}
\usepackage{booktabs}
\usepackage{array}
\usepackage{caption}

\usepackage[numbers,sort&compress]{natbib}
\usepackage{xcolor}
\usepackage{enumitem}
\usepackage{authblk}

\definecolor{worknotebg}{RGB}{245,247,250}
\definecolor{worknotefg}{RGB}{70,78,90}
\definecolor{authorqueryborder}{RGB}{150,35,35}
\definecolor{authorquerybg}{RGB}{255,250,230}
\newif\ifworkingdraft
\workingdrafttrue
\newcommand{\worknote}[1]{%
  \ifworkingdraft
    \par\smallskip
    \noindent\colorbox{worknotebg}{%
      \parbox{0.96\linewidth}{\small\color{worknotefg}\textbf{Робоча примітка.} #1}}%
    \par\smallskip
  \fi}
\newcommand{\AuthorQuery}[2]{%
  \ifworkingdraft
    \par\smallskip\noindent
    {\footnotesize
    \fcolorbox{authorqueryborder}{authorquerybg}{%
      \parbox{\dimexpr\linewidth-2\fboxsep-2\fboxrule\relax}{%
        \textcolor{authorqueryborder}{\textbf{Питання автору #1.}}%
        \textcolor{black}{\ #2}}}%
    }%
    \par\smallskip
  \fi}

\newcommand{\resultfigure}[4][0.98\linewidth]{%
  \begin{figure}[htbp]
    \centering
    \IfFileExists{../../figures/#2}{%
      \includegraphics[width=#1]{#2}%
    }{%
      \fbox{\parbox[c][45mm][c]{0.9\linewidth}{\centering
        Файл рисунка ще не знайдено: \texttt{\detokenize{#2}}}}%
    }
    \caption{#3}
    \label{#4}
  \end{figure}}

\newcommand{\dd}{\mathop{}\!\mathrm{d}}
\newcommand{\ii}{\mathrm{i}}
\newcommand{\ee}{\mathrm{e}}
\newcommand{\sigmap}{\sigma^{\prime}}
\newcommand{\deltap}{\delta^{\prime}}
\newcommand{\Jp}{J^{\prime}}
\newcommand{\jump}[1]{\left[\!\left[#1\right]\!\right]}

\usepackage[hidelinks]{hyperref}
\hypersetup{
  pdftitle={Відгалуження міжфазної зсувної тріщини у кутовій точці ламаної межі поділу},
  pdfauthor={А. О. Камінський; М. Ф. Селіванов; М. В. Дудик; Т. В. Поліщук}
}
\setlength{\parindent}{1.25em}
\setlength{\parskip}{0pt}
\setlength{\emergencystretch}{4em}
\captionsetup{font=small,labelfont=bf}

\newif\ifbibliographyready
\bibliographyreadytrue
